% GENERATED FILE. Edit canonical lessons, then run npm run generate:books.
% canonical-source-hash: fnv1a64:4d6f94976fbd6131
% canonical-lessons: KA-C384-grahisu, KA-C384-vyakhyana, KA-C384-kalpane, KA-C384-amsa, KA-C384-laksana

\chapter{Grasping an Idea}
\label{ch:ka-grasping-an-idea}
\hlchaptermodality{384}

\begin{chapteropening}
\textbf{By the end of this chapter:} \emph{I can say to grasp, an interpretation, a notion, an aspect and a feature.}

Complete the last lesson of chapter 384: I can say to grasp, an interpretation, a notion, an aspect and a feature.
\end{chapteropening}

\section[\texorpdfstring{\kn{ಗ್ರಹಿಸು} (grahisu)}{ಗ್ರಹಿಸು (grahisu)}]{\kn{ಗ್ರಹಿಸು} (grahisu) — to grasp; to comprehend}

\label{lesson:KA-C384-grahisu}



\begin{quote}
Before the new one: say the Kannada for impossible, then the Kannada for to deny.
\end{quote}

\subsection*{You'll want to know: \kn{ಗ್ರಹಿಸು}}
\textbf{\kn{ಗ್ರಹಿಸು}} — \emph{grahisu} — \textquotedblleft{}to grasp; to comprehend\textquotedblright{}.

\begin{quote}
\kn{ಮಗು} \kn{ಪಾಠವನ್ನು} \kn{ಬೇಗ} \kn{ಗ್ರಹಿಸಿತು}.
\end{quote}

\begin{quote}
\emph{magu pāṭhavannu bēga grahisitu.}
\end{quote}

\begin{quote}
The child grasped the lesson quickly.
\end{quote}

Past (I …): \kn{ಗ್ರಹಿಸಿದೆ}, \emph{grahiside}.

\kn{ಗ್ರಹಿಕೆ}, understanding.

\begin{cousinweb}
\textbf{\kn{ಗ್ರಹಿಸು}} (\emph{grahisu}) is built on Sanskrit \emph{grahaṇa}, \textquotedblleft{}seizing, grasping\textquotedblright{} (CDIAL 4364), with the Kannada verb ending -\emph{isu}.
\end{cousinweb}

\begin{grammarlens}[title={What you've built}]
One more word for giving reasons and opinions.
\end{grammarlens}

\subsection*{Your turn}
\begin{itemize}
\raggedright
  \item \emph{Say it:} \emph{grahisu}
  \item \emph{Say it:} \emph{grahisu}, once more
  \item \emph{Say it:} \emph{magu pāṭhavannu bēga grahisitu.} — The child grasped the lesson quickly.
  \item \emph{Recall:} say the Kannada for impossible, then the Kannada for to deny, then say \emph{grahisu} again
\end{itemize}

\begin{tcolorbox}[breakable,colback=teal!4,colframe=teal!35!black,title={Before you move on}]
What does \kn{ಗ್ರಹಿಸು} mean? (\textquotedblleft{}To grasp; to comprehend\textquotedblright{}.) Say it once more.
\end{tcolorbox}
\section[\texorpdfstring{\kn{ವ್ಯಾಖ್ಯಾನ} (vyākhyāna)}{ವ್ಯಾಖ್ಯಾನ (vyākhyāna)}]{\kn{ವ್ಯಾಖ್ಯಾನ} (vyākhyāna) — an interpretation}

\label{lesson:KA-C384-vyakhyana}



\begin{quote}
Before the new one: say the Kannada for to deny, then the Kannada for to grasp.
\end{quote}

\subsection*{You'll want to know: \kn{ವ್ಯಾಖ್ಯಾನ}}
\textbf{\kn{ವ್ಯಾಖ್ಯಾನ}} — \emph{vyākhyāna} — \textquotedblleft{}an interpretation\textquotedblright{}.

\begin{quote}
\kn{ಆ} \kn{ಹಳೆಯ} \kn{ಹಾಡಿಗೆ} \kn{ಹಲವು} \kn{ವ್ಯಾಖ್ಯಾನಗಳು} \kn{ಇವೆ}.
\end{quote}

\begin{quote}
\emph{ā haḷeya hāḍige halavu vyākhyānagaḷu ive.}
\end{quote}

\begin{quote}
That old song has many interpretations.
\end{quote}

Plural: \kn{ವ್ಯಾಖ್ಯಾನಗಳು}, \emph{vyākhyānagaḷu}.

Also a commentary on a text.

\begin{cousinweb}
\textbf{\kn{ವ್ಯಾಖ್ಯಾನ}} (\emph{vyākhyāna}) is borrowed from Sanskrit \emph{vyākhyāna}, \textquotedblleft{}explaining, commenting\textquotedblright{} (CDIAL 12188).
\end{cousinweb}

\begin{grammarlens}[title={What you've built}]
One more word for giving reasons and opinions.
\end{grammarlens}

\subsection*{Your turn}
\begin{itemize}
\raggedright
  \item \emph{Say it:} \emph{vyākhyāna}
  \item \emph{Say it:} \emph{vyākhyāna}, once more
  \item \emph{Say it:} \emph{ā haḷeya hāḍige halavu vyākhyānagaḷu ive.} — That old song has many interpretations.
  \item \emph{Recall:} say the Kannada for to deny, then the Kannada for to grasp, then say \emph{vyākhyāna} again
\end{itemize}

\begin{tcolorbox}[breakable,colback=teal!4,colframe=teal!35!black,title={Before you move on}]
What does \kn{ವ್ಯಾಖ್ಯಾನ} mean? (\textquotedblleft{}An interpretation\textquotedblright{}.) Say it once more.
\end{tcolorbox}
\section[\texorpdfstring{\kn{ಕಲ್ಪನೆ} (kalpane)}{ಕಲ್ಪನೆ (kalpane)}]{\kn{ಕಲ್ಪನೆ} (kalpane) — a notion; an idea}

\label{lesson:KA-C384-kalpane}



\begin{quote}
Before the new one: say the Kannada for to grasp, then the Kannada for an interpretation.
\end{quote}

\subsection*{You'll want to know: \kn{ಕಲ್ಪನೆ}}
\textbf{\kn{ಕಲ್ಪನೆ}} — \emph{kalpane} — \textquotedblleft{}a notion; an idea\textquotedblright{}.

\begin{quote}
\kn{ಇದು} \kn{ಒಳ್ಳೆಯ} \kn{ಕಲ್ಪನೆ}.
\end{quote}

\begin{quote}
\emph{idu oḷḷeya kalpane.}
\end{quote}

\begin{quote}
This is a good idea.
\end{quote}

Also imagination: \kn{ಕಲ್ಪನೆಯ} \kn{ಕಥೆ}, a made-up story.

\begin{cousinweb}
\textbf{\kn{ಕಲ್ಪನೆ}} (\emph{kalpane}) is borrowed from Sanskrit \emph{kalpanā}, \textquotedblleft{}making, forming, arranging\textquotedblright{} (Monier-Williams), from \emph{kḷp}, \textquotedblleft{}to be fit, to arrange\textquotedblright{}.
\end{cousinweb}

\begin{grammarlens}[title={What you've built}]
One more word for giving reasons and opinions.
\end{grammarlens}

\subsection*{Your turn}
\begin{itemize}
\raggedright
  \item \emph{Say it:} \emph{kalpane}
  \item \emph{Say it:} \emph{kalpane}, once more
  \item \emph{Say it:} \emph{idu oḷḷeya kalpane.} — This is a good idea.
  \item \emph{Recall:} say the Kannada for to grasp, then the Kannada for an interpretation, then say \emph{kalpane} again
\end{itemize}

\begin{tcolorbox}[breakable,colback=teal!4,colframe=teal!35!black,title={Before you move on}]
What does \kn{ಕಲ್ಪನೆ} mean? (\textquotedblleft{}A notion; an idea\textquotedblright{}.) Say it once more.
\end{tcolorbox}
\section[\texorpdfstring{\kn{ಅಂಶ} (aṃśa)}{ಅಂಶ (aṃśa)}]{\kn{ಅಂಶ} (aṃśa) — an aspect; a point}

\label{lesson:KA-C384-amsa}



\begin{quote}
Before the new one: say the Kannada for an interpretation, then the Kannada for a notion.
\end{quote}

\subsection*{You'll want to know: \kn{ಅಂಶ}}
\textbf{\kn{ಅಂಶ}} — \emph{aṃśa} — \textquotedblleft{}an aspect; a point\textquotedblright{}.

\begin{quote}
\kn{ಈ} \kn{ಪಾಠದಲ್ಲಿ} \kn{ಮೂರು} \kn{ಮುಖ್ಯ} \kn{ಅಂಶಗಳು} \kn{ಇವೆ}.
\end{quote}

\begin{quote}
\emph{ī pāṭhadalli mūru mukhya aṃśagaḷu ive.}
\end{quote}

\begin{quote}
This lesson has three main points.
\end{quote}

Plural: \kn{ಅಂಶಗಳು}, \emph{aṃśagaḷu}.

Also a share or a part.

\begin{cousinweb}
\textbf{\kn{ಅಂಶ}} (\emph{aṃśa}) is borrowed from Sanskrit \emph{aṃśa}, \textquotedblleft{}a share, a portion, a part\textquotedblright{} (Monier-Williams).
\end{cousinweb}

\begin{grammarlens}[title={What you've built}]
One more word for giving reasons and opinions.
\end{grammarlens}

\subsection*{Your turn}
\begin{itemize}
\raggedright
  \item \emph{Say it:} \emph{aṃśa}
  \item \emph{Say it:} \emph{aṃśa}, once more
  \item \emph{Say it:} \emph{ī pāṭhadalli mūru mukhya aṃśagaḷu ive.} — This lesson has three main points.
  \item \emph{Recall:} say the Kannada for an interpretation, then the Kannada for a notion, then say \emph{aṃśa} again
\end{itemize}

\begin{tcolorbox}[breakable,colback=teal!4,colframe=teal!35!black,title={Before you move on}]
What does \kn{ಅಂಶ} mean? (\textquotedblleft{}An aspect; a point\textquotedblright{}.) Say it once more.
\end{tcolorbox}
\section[\texorpdfstring{\kn{ಲಕ್ಷಣ} (lakṣaṇa)}{ಲಕ್ಷಣ (lakṣaṇa)}]{\kn{ಲಕ್ಷಣ} (lakṣaṇa) — a feature; a mark}

\label{lesson:KA-C384-laksana}



\begin{quote}
Before the new one: say the Kannada for a notion, then the Kannada for an aspect.
\end{quote}

\subsection*{You'll want to know: \kn{ಲಕ್ಷಣ}}
\textbf{\kn{ಲಕ್ಷಣ}} — \emph{lakṣaṇa} — \textquotedblleft{}a feature; a mark\textquotedblright{}.

\begin{quote}
\kn{ತಾಳ್ಮೆ} \kn{ಒಳ್ಳೆಯ} \kn{ಶಿಕ್ಷಕನ} \kn{ಲಕ್ಷಣ}.
\end{quote}

\begin{quote}
\emph{tāḷme oḷḷeya śikṣakana lakṣaṇa.}
\end{quote}

\begin{quote}
Patience is the mark of a good teacher.
\end{quote}

\kn{ಲಕ್ಷಣವಾಗಿ}, neatly, handsomely.

\begin{cousinweb}
\textbf{\kn{ಲಕ್ಷಣ}} (\emph{lakṣaṇa}) is borrowed from Sanskrit \emph{lakṣaṇa}, \textquotedblleft{}a mark, a sign, a characteristic\textquotedblright{} (CDIAL 10882).
\end{cousinweb}

\begin{grammarlens}[title={What you've built}]
One more word for giving reasons and opinions.
\end{grammarlens}

\subsection*{Your turn}
\begin{itemize}
\raggedright
  \item \emph{Say it:} \emph{lakṣaṇa}
  \item \emph{Say it:} \emph{lakṣaṇa}, once more
  \item \emph{Say it:} \emph{tāḷme oḷḷeya śikṣakana lakṣaṇa.} — Patience is the mark of a good teacher.
  \item \emph{Recall:} say the Kannada for a notion, then the Kannada for an aspect, then say \emph{lakṣaṇa} again
\end{itemize}

\begin{tcolorbox}[breakable,colback=teal!4,colframe=teal!35!black,title={Before you move on}]
What does \kn{ಲಕ್ಷಣ} mean? (\textquotedblleft{}A feature; a mark\textquotedblright{}.) Say it once more.
\end{tcolorbox}
